\section{Introduction}
\label{sec:introduction}
% Role: opening; define the workload and the systems question.
Panorama stitching combines overlapping images into a wider field of view by
estimating their alignment and composing the warped image
content~\cite{brown2007stitching}. A GPU implementation parallelizes the
operations but does not determine how much computation an input pair needs. This work
studies budget allocation within a feature-based pipeline that produces a
panorama from two original RGB images.

% Role: challenge; state a testable hypothesis rather than an empirical result.
The computational budgets of this pipeline interact across stages. Reducing the
alignment resolution changes the detected keypoints and descriptors, limiting
the keypoints changes the matching workload, and the robust-fitting budget
determines which homography hypotheses are considered. A cheaper choice can also
produce an unusable alignment. The relevant objective is therefore latency under
an explicit quality policy rather than the execution time of an isolated fast
kernel, and this latency includes the cost of unsuccessful attempts.

% Role: method scope; distinguish implemented components from pending runtime work.
We expose a configuration $c=(s,k,r,b)$ comprising alignment-image scale,
keypoint limit, robust-fitting hypothesis budget, and blending mode. Rendering
uses the original RGB images and a homography in their coordinate systems.
The proposed runtime selects a configuration from information available at
decision time and can increase budgets after inspecting intermediate signals.
Selection, checks, synchronization, retries, and reuse bookkeeping must be
included in end-to-end latency. An offline oracle that already knows which
configurations satisfy the quality policy is a reference bound, not an online
selection algorithm.

% Role: research question and evidence obligations.
The research question is: \emph{Can a scene-aware cross-stage runtime reduce GPU
panorama-stitching latency by dynamically allocating computational budgets while
maintaining a predefined quality-of-result constraint?} The implementation
provides a CUDA pipeline and rendering policies with explicit geometry, sampling,
and buffer contracts. The held-out T4 comparison records latency, relative
alignment quality and explicit failure dispositions. The frozen heuristic
provides a conditional latency benefit on common passing pairs, accompanied
by reduced quality coverage. These measurements bound the claim supported
by that controller.

% Role: positioning and organization; avoid unsupported novelty claims.
The intended systems contribution is an evaluated policy for coordinating
budgets and valid result reuse, rather than a new stitching or blending algorithm.
Section~\ref{sec:related} positions the study against stitching, composition,
and configuration tuning. Section~\ref{sec:pipeline} describes the checked pipeline
interface and renderer. Section~\ref{sec:controller} describes the controller.
Section~\ref{sec:experiments} evaluates the frozen T4 configurations, and
Section~\ref{sec:discussion} examines their limitations.
